\begin{RequAnswer}{2.1}
Answers will vary, but it should say something like
``An argument using logic and/or math to demonstrate or show the nature of a conclusion,''
or
``the process or an instance of establishing the validity of a statement especially by derivation
from other statements in accordance with principles of reasoning'' (the latter is from Merriam-Webster).
\end{RequAnswer}
\begin{RequAnswer}{2.2}
If someone correctly proves statement $A$, that means it is a true statement, regardless of whether or not
you understand the proof.  That's because, by definition, a proof establishes the validity of a statement.
\end{RequAnswer}
\begin{RequAnswer}{2.3}
True. An even integer is one of the form $2k$, where $k$ is an integer.
An odd integer is one of the form $2k + 1$ where $k$ is an integer.
Thus, even numbers are divisible by 2 whereas odd numbers are not.
A number cannot both be divisible by 2 and not divisible by 2 at the same time.
\end{RequAnswer}
\begin{RequAnswer}{2.4}
No. For instance, 6 is divisible by 2, but 2 is not divisible by 6.
\end{RequAnswer}
\begin{RequAnswer}{2.5}
No.
A number cannot be both composite and prime because by definition,
a composite integer is a positive integer $c>1$ that is {\em not} prime.
Thus, every integer greater than $1$ is either prime or composite, but never both.
\end{RequAnswer}
\begin{RequAnswer}{2.6}
3, 97, 173, and 999983 are prime.
4, 6, 27, 38, 150, and 999985 are composite.
1 is neither. It is impossible to be both.
\end{RequAnswer}
\begin{RequAnswer}{2.7}
$6!=6\cdot 5\cdot 4\cdot 3\cdot 2\cdot 1 = 720$.
$7!=7\cdot 6! = 7*720 =5040,$ computed the easy way
be recognizing I just needed to multiply the previous answer by 7.
\end{RequAnswer}
\begin{RequAnswer}{2.8}
Clearly $2!=2$ is prime.
We will show that this is the only case in which $n!$ is prime.
$1!=1$, which is not prime by definition.
$3!=2\cdot 3$, which is clearly not prime.
If $n>3$, $n!$ is divisible by $3!=3\cdot 2$, so it not prime.
\end{RequAnswer}
\begin{RequAnswer}{2.9}
No.
Take the example from the book, ``If you know Java, then you know a programming language'' (where $A$ is the proposition
``you know Java'' and $B$ is the proposition ``you know a programming language.''
Since Java is a programming language, this proposition is true.
Then $B$ implies $A$ is the proposition ``if you know a programming language, then you know Java.'' But that is clearly not
true. You might know {\em C++}, which is a programming language, but not know Java.
\end{RequAnswer}
\begin{RequAnswer}{2.10}
False. From the previous answer, consider the implication ``if you know a programming language, then you know Java,''
and its inverse, ``if you do not know a programming language, then you do not know Java."
It shouldn't be too difficult to see that although the inverse is true, the implication is false.
\end{RequAnswer}
\begin{RequAnswer}{2.11}
This one is true since the inverse and converse of an implication are contrapositives of each other, and an implication
and its contrapositive are equivalent.
\end{RequAnswer}
\begin{RequAnswer}{2.12}
An implication and its contrapositive are equivalent.
Therefore, if one is true, then the other is true
(and if one is false, the other is false, of course).
\end{RequAnswer}
\begin{RequAnswer}{2.13}
Your answer will likely be different, and hopefully more detailed than the one provided here.
But you should say something about how you assume that what you want to prove is false,
then use logic to arrive at a contradiction (that is, a statement that you know to be false).
Since you ``proved'' a false statement using correct logic, it must be that your premise (that is,
the statement that you assumed was false) is incorrect. Since your premise was that the statement
you wanted to prove was false, then it must be that the statement is true (again, because you
showed that if it is false, then you can prove something that is not true).
\end{RequAnswer}
\begin{RequAnswer}{2.14}
They are
(cow, chicken, rabbit),
(cow, rabbit, chicken),
(chicken, cow, rabbit),
(chicken, rabbit, cow),
(rabbit, cow, chicken), and
(rabbit, chicken, cow).
\end{RequAnswer}
\begin{RequAnswer}{2.15}
True. Every integer $a$ can be written as $a=\frac{a}{1}$, where $a$ and $1$ are both integers and $1\neq 0$, so it is rational.
\end{RequAnswer}
\begin{RequAnswer}{2.16}
Assume that $k$ is the smallest positive rational number.
Since $k$ is rational, we can write it as $p/q$ for integers
$p$ and $q\neq 0$. But clearly $k/2=p/(2q)$ is positive,
rational, and smaller than $k$, which contradicts our assumption
that $k$ was the smallest such number.
Therefore there is no smallest positive rational number.
\end{RequAnswer}
\begin{RequAnswer}{2.17}
Since an implication and its contrapositive are equivalent, if you prove the contrapositive of a implication is
true, then the implication must also be true. And that is exactly what proof by contraposition does.
\end{RequAnswer}
\begin{RequAnswer}{2.18}
{\bf Contradiction:} Assume that $p$ and $\neg q$ are both true. Get a contradiction.
Conclude that if $p$ is true, $\neg q$ cannot be true, so $q$ must be true. Thus, $p\rightarrow q$ is true.

\noindent
{\bf Contraposition: } Prove that $\neg q\rightarrow \neg p$, which is equivalent to $p\rightarrow q$.

\noindent
In both cases, you assume that $\neg q$ is true. Often the contradiction you get is that $\neg p$
is true (which is the goal in contraposition proof),
so the majority of the proofs look the same. But they begin and end slightly differently.
\end{RequAnswer}
\begin{RequAnswer}{2.19}
True. From the Question~\ref{intIsRat}, we know that every integer is rational, so integers are not irrational.
Therefore, if a number is irrational, it cannot be an integer.
Alternatively, you can recognize that this is essentially the contrapositive of the Question~\ref{intIsRat}, so it is also true.
\end{RequAnswer}
\begin{RequAnswer}{2.20}
(a) Assume that $x>0$ is irrational, but that $\sqrt{x}$ is rational.
Then $\sqrt{x}=p/q$ for some integers $p, q\neq 0$.
That means that $x=(\sqrt{x})^2=(p/q)^2=p^2/q^2$, which
is clearly rational since $p^2$ and $q^2\neq 0$ are both rational.
But this contradicts our assumption that $x$ is irrational.
Therefore, $\sqrt{x}$ must be irrational.\\
(b) We will prove that if $\sqrt{x}$ is rational, then $x$ is rational,
which will imply our statement is true since this is the
contrapositive of our statement.
Assume $\sqrt{x}$ is rational. Then
$\sqrt{x}=p/q$ for some integers $p, q\neq 0$.
Therefore, $x=(\sqrt{x})^2=(p/q)^2=p^2/q^2$
which is clearly rational since $p^2$ and $q^2\neq 0$ are both rational.
\end{RequAnswer}
\begin{RequAnswer}{2.21}
False. For instance, $1.5=\frac{3}{2}$ is rational, but clearly not an integer.
\end{RequAnswer}
\begin{RequAnswer}{2.22}
We will use a proof by cases.
Case 1: If $n$ is even, then $n=2k$ for some integer $k$.
Then $n^2=(2k)^2=4k^2=2(2k^2)$, which is even since
$2k^2$ is an integer.
Case 2: If $n$ is odd, then $n=2k+1$ for some integer $k$.
Then $n^2=(2k+1)^2=4k^2+4k+1=2(2k^2+2k)+1$, which is odd since $2k^2+2k$ is an integer.
In both cases, the parity remains unchanged.
\end{RequAnswer}
\begin{RequAnswer}{2.23}
No you would have to show the reverse as well.
That is, if I want to prove $A$ if and only if $B$,
I would have to prove either the proposition ($A\rightarrow B$) or the contrapositive ($\neg B \rightarrow \neg A$)
{\bf and} I would have to prove the inverse ($\neg A\rightarrow \neg B$) or the converse ($B\rightarrow A$).
\end{RequAnswer}
\begin{RequAnswer}{2.24}
(a) No. This is proving the forward direction twice by proving the implication and its contrapositive.
One of these needs to be replaced with either the inverse or converse.
(b) Yes. This is proving the contrapositive and the converse.
\end{RequAnswer}
\begin{RequAnswer}{2.25}
For the forward direction, we assume that $n$ is even.
Then $n=2k$ for some integer $k$,
so $n^2=(2k)^2=4k^2=2(2k^2)$, which is even since
$2k^2$ is an integer. Thus, $n$ even implies $n^2$ is even.\\
For the reverse direction, we will prove that
$n$ is not even implies $n^2$ is not even. In other words,
$n$ is odd implies that $n^2$ is odd.
Assume $n$ is odd. Then $n=2k+1$ for some integer $k$,
so $n^2=(2k+1)^2=4k^2+4k+1=2(2k^2+2k)+1$, which is odd since $2k^2+2k$ is an integer.
Thus, $n$ odd implies that $n^2$ is odd, which (as we have already
said, but we'll say it again anyway) is equivalent to
$n^2$ is even implies $n$ is even.
\end{RequAnswer}
\begin{RequAnswer}{2.26}
You use proof by counterexample to {\em disprove} statements. You only need to show
one instance where the statement is not true to demonstrate that it is not always true, so proof by counterexample
is valid. On the other hand, proof by example is just demonstrating the truth of a statement given some specific values.
Just because it works for the given values, that does not prove that it works for all values, so it is not a valid proof technique.
\end{RequAnswer}
\begin{RequAnswer}{2.27}
This is definitely {\em not} a valid proof technique! The problem is that when you write down an equation and start
working both sides of it, you are implicitly assuming that the equation you are starting with is true. But since you
are trying to prove that it is true, you can't start your proof with the fact that it is true--that is circular reasoning.
For instance, do you remember the supposed proof of  $-1=1$ that started by working both sides of the equation?
\end{RequAnswer}
\begin{RequAnswer}{2.28}
The step that divided both sides by $a-b$ was division by zero,
which is not allowed!
\end{RequAnswer}
